\documentclass[11pt]{article}

\usepackage[margin=1in]{geometry}
\usepackage{amsmath,amssymb,amsthm}
\usepackage{enumitem}

\newtheorem{theorem}{Theorem}
\newtheorem{proposition}{Proposition}

\title{A Lehmann--Romano Proof of Admissibility for Closed Convex Acceptance Regions}
\date{}

\begin{document}
\maketitle

\section{Lehmann--Romano supported formulation}

The closed-convex acceptance-region argument in Lehmann and Romano, Section 6.7, gives a sufficiency result, not an unrestricted if-and-only-if result.  The conclusion has two parts.  First, a nonrandomized test with closed convex acceptance region is $d$-admissible, once the ray condition in Theorem 6.7.1 is verified.  Second, if the same test has size exactly $\alpha$ at a finite null point, then Corollary 6.7.1 gives $\alpha$-admissibility.

Thus the result proved below is
\[
A_0\text{ closed and convex}\quad\Longrightarrow\quad d\text{-admissibility},
\]
and, with exact size attained at a finite point of the null,
\[
A_0\text{ closed and convex}\quad\Longrightarrow\quad \alpha\text{-admissibility}.
\]
The exact-size condition is needed for the $\alpha$-admissibility conclusion, but not for the $d$-admissibility conclusion.

\section{The theorem and proof}

\begin{theorem}[Closed convex acceptance regions imply admissibility]
Let $0<\alpha<1$. Let
\[
X\sim N(\mu,\Sigma),
\]
where $\Sigma$ is known and positive definite and $\mu\in\mathbf R^k$. Consider testing
\[
H_0:\mu\in\Lambda_0=\{\mu:\mu\le 0\}
\quad\text{versus}\quad
H_1:\mu\in\Lambda_1=\mathbf R^k\setminus\Lambda_0.
\]
Let $\phi_0$ be a nonrandomized test with acceptance region
\[
A_0=\{x:\phi_0(x)=0\}.
\]
Assume:
\begin{enumerate}[label=(\roman*)]
\item $A_0$ is closed and convex;
\item $\phi_0$ has level $\alpha$:
\[
\sup_{\mu\in\Lambda_0} E_\mu[\phi_0(X)]\le \alpha.
\]
\end{enumerate}
Then $\phi_0$ is $d$-admissible for testing $H_0$ against $H_1$. If, in addition, there exists a finite point $\mu_0\in\Lambda_0$ such that
\[
E_{\mu_0}[\phi_0(X)]=\alpha,
\]
then $\phi_0$ is $\alpha$-admissible.
\end{theorem}

\begin{proof}
We use Theorem 6.7.1 and Corollary 6.7.1 of Lehmann and Romano.  The relevant result is the following.

\begin{proposition}[Lehmann--Romano]
Consider an exponential family with density
\[
p_\theta(t)=C(\theta)\exp\{\theta't\}
\]
with respect to a $\sigma$-finite measure $\nu$, with natural parameter space $\Omega$. Let $\Omega_H$ and $\Omega'$ be disjoint nonempty subsets of $\Omega$. Let $\phi_0$ be a nonrandomized test of $H:\theta\in\Omega_H$, based on $T$, with acceptance region $A_0\subseteq\mathbf R^s$. Suppose $A_0$ is closed and convex and that the following ray condition holds: whenever
\[
A_0\cap\{t:a't>c\}=\varnothing,
\]
there exist $\theta^*\in\Omega$ and $\lambda_n\to\infty$ such that
\[
\theta^*+\lambda_n a\in\Omega'.
\]
Then any competing acceptance region $A$ satisfying
\[
P_\theta(T\in A)\le P_\theta(T\in A_0)
\qquad\text{for all }\theta\in\Omega'
\]
is contained in $A_0$, except for a $\nu$-null set. Consequently, the test with acceptance region $A_0$ is $d$-admissible. If the test has size $\alpha$ and there exists a finite point $\theta_0\in\overline{\Omega_H}$ such that
\[
E_{\theta_0}[\phi_0(X)]=\alpha,
\]
then the test is also $\alpha$-admissible. The same conclusions allow randomized competing tests.
\end{proposition}

We now verify the assumptions of this result.  Write the normal density in canonical exponential-family form. Put
\[
\theta=\Sigma^{-1}\mu,
\qquad
\mu=\Sigma\theta.
\]
With respect to the measure
\[
d\nu(x)=(2\pi)^{-k/2}|\Sigma|^{-1/2}\exp\left\{-\frac12 x'\Sigma^{-1}x\right\}dx,
\]
the density of $X\sim N(\Sigma\theta,\Sigma)$ is
\[
p_\theta(x)=\exp\left\{\theta'x-\frac12\theta'\Sigma\theta\right\}.
\]
Thus $T(x)=x$, $\Omega=\mathbf R^k$, and
\[
\Omega_H=\{\theta:\Sigma\theta\le 0\},
\qquad
\Omega_K=\{\theta:\Sigma\theta\not\le 0\}.
\]
We apply the Lehmann--Romano result with $\Omega'=\Omega_K$.

It remains to verify the ray condition.  Let $a\ne 0$ and $c\in\mathbf R$ satisfy
\[
A_0\cap\{x:a'x>c\}=\varnothing.
\]
Equivalently, $A_0\subseteq\{x:a'x\le c\}$.  We claim that
\[
\Sigma a\not\le 0.
\]
Suppose instead that $\Sigma a\le 0$.  For $\lambda>0$, define
\[
\theta_\lambda=\lambda a,
\qquad
\mu_\lambda=\Sigma\theta_\lambda=\lambda\Sigma a.
\]
Then $\mu_\lambda\le 0$, so $\theta_\lambda\in\Omega_H$.  Under $P_{\theta_\lambda}$,
\[
a'X\sim N\left(\lambda a'\Sigma a,\ a'\Sigma a\right).
\]
Since $\Sigma$ is positive definite and $a\ne 0$, $a'\Sigma a>0$. Hence
\[
P_{\theta_\lambda}(a'X>c)\to 1.
\]
Because $A_0\subseteq\{x:a'x\le c\}$, we have $\{x:a'x>c\}\subseteq A_0^c$. Since $\phi_0$ is nonrandomized,
\[
E_{\theta_\lambda}[\phi_0(X)]
\ge P_{\theta_\lambda}(a'X>c)
\to 1.
\]
This contradicts the level condition because every $\theta_\lambda\in\Omega_H$ and $\alpha<1$. Therefore $\Sigma a\not\le 0$.

It follows that $\lambda\Sigma a\not\le 0$ for every $\lambda>0$. Hence
\[
\lambda a\in\Omega_K
\qquad\text{for every }\lambda>0.
\]
The ray condition holds with $\theta^*=0$ and any sequence $\lambda_n\to\infty$.

The Lehmann--Romano proposition now implies $d$-admissibility of $\phi_0$.

Finally suppose there exists finite $\mu_0\in\Lambda_0$ such that $E_{\mu_0}[\phi_0(X)]=\alpha$.  Let $\theta_0=\Sigma^{-1}\mu_0$.  Then $\theta_0\in\Omega_H\subseteq\overline{\Omega_H}$ and
\[
E_{\theta_0}[\phi_0(X)]=\alpha.
\]
The second part of the Lehmann--Romano corollary therefore gives $\alpha$-admissibility.
\end{proof}

\section{Why exact size is needed for $\alpha$-admissibility}

The statement ``closed and convex acceptance region implies $\alpha$-admissibility'' is false if one only assumes
\[
\sup_{\mu\le 0}E_\mu[\phi(X)]\le \alpha.
\]
For example, the zero test $\phi(x)\equiv 0$ has acceptance region $\mathbf R^k$, which is closed and convex, and it satisfies the size constraint for every $0<\alpha<1$. But it is not $\alpha$-admissible. The test
\[
\psi(x)=1\{x_1>\sqrt{\Sigma_{11}}\Phi^{-1}(1-\alpha)\}
\]
has level $\alpha$ for $H_0:\mu\le 0$ and has strictly positive power at alternatives with $\mu_1>0$, while the zero test has zero power everywhere.

This example does not contradict $d$-admissibility. A $d$-dominating test would also have to satisfy
\[
E_\mu[\psi(X)]\le E_\mu[\phi(X)]=0
\qquad\text{for all }\mu\in\Lambda_0,
\]
which forces no increase in rejection probability on the null. Thus exact size is needed for the $\alpha$-admissibility conclusion, not for the $d$-admissibility conclusion.

\begin{thebibliography}{9}

\bibitem{LehmannRomano2005}
Lehmann, E. L. and Romano, J. P. (2005).
\textit{Testing Statistical Hypotheses}.
Third edition. Springer.

\end{thebibliography}

\end{document}
